%% ma: L11-B14-0016
%% lop: 11
%% chuong: 4
%% bai: 14
%% dang: 11.14.3
%% loai: TN4
%% muc: NB
%% dap_an: "A"
%% nguon: "tai-lieu-hoc-tap-mon-toan-lop-11-theo-chuong-trinh-giao-duc-pho-thong-2018.pdf (Lớp Toán Thầy Duy, Phạm Lê Duy, THPT Chu Văn An), Chương 4 Bài 5, Phần 1, Câu 11 (tr.529)"
%% kien_thuc: "Lớp 11 Bài 14 (phép chiếu song song)"
%% trang_thai: chinh-thuc
%% ghi_chu: "Đã duyệt 03/11/2026."
\begin{ex}
Hình biểu diễn của hình chữ nhật trong không gian \textbf{không} thể là hình nào trong các hình sau?
\choice{\True Hình thang.}{Hình bình hành.}{Hình chữ nhật.}{Hình thoi.}
\loigiai{\sloppy Hình biểu diễn của hình chữ nhật là hình bình hành (có thể là hình chữ nhật hoặc hình thoi), vì phép chiếu song song bảo toàn tính song song của các cặp cạnh đối. Hình thang (có hai cạnh không song song) không thể là hình biểu diễn.}
\end{ex}
